% Shared notation for the differential-geometric SN Ia light-curve project.
% Input this file from every scientific document so notation stays consistent.
%
% Requires: amsmath, amssymb

% --- Number sets -------------------------------------------------------
\newcommand{\Real}{\mathbb{R}}
\newcommand{\Rn}{\mathbb{R}^{n}}
\newcommand{\Rtwo}{\mathbb{R}^{2}}
\newcommand{\Rthree}{\mathbb{R}^{3}}

% --- Differential geometry ---------------------------------------------
\newcommand{\dd}{\mathrm{d}}
\newcommand{\curve}{\gamma}                 % the light-curve curve
\newcommand{\arc}{s}                        % arclength
\newcommand{\curv}{\kappa}                  % curvature
\newcommand{\tors}{\tau}                    % torsion (n >= 3 only)
\newcommand{\frenet}[1]{\mathbf{E}_{#1}}    % Frenet frame vector
\newcommand{\tang}{\mathbf{T}}              % unit tangent
\newcommand{\norml}{\mathbf{N}}             % unit normal
\newcommand{\arclen}{L}                     % total arclength of the window

% --- Model parameters ---------------------------------------------------
\newcommand{\shape}{\theta}                 % shape latents
\newcommand{\shapeone}{\theta_{1}}
\newcommand{\shapetwo}{\theta_{2}}
\newcommand{\tscale}{\sigma}               % timing scale
\newcommand{\tmax}{t_{\mathrm{max}}}        % time origin
\newcommand{\warpone}{a_{1}}                % quadratic timing corrections
\newcommand{\warptwo}{a_{2}}
\newcommand{\dm}{\mu}                       % normalization / distance modulus
\newcommand{\colr}{c}                       % colour translation (variant only)

% --- Photometry ---------------------------------------------------------
\newcommand{\magg}{m_{g}}
\newcommand{\magr}{m_{r}}
\newcommand{\fluxg}{f_{g}}
\newcommand{\fluxr}{f_{r}}
\newcommand{\zp}{\mathrm{zp}}
\newcommand{\phase}{p}                      % rest-frame phase
\newcommand{\zred}{z}

% --- Objects and abbreviations -----------------------------------------
\newcommand{\snia}{SN\,Ia}
\newcommand{\sneia}{SNe\,Ia}
\newcommand{\ztf}{ZTF}
\newcommand{\salt}{SALT2}
\newcommand{\NN}{\mathrm{NN}}

% --- Operators ----------------------------------------------------------
\DeclareMathOperator{\SO}{SO}
\DeclareMathOperator*{\argmin}{arg\,min}
\newcommand{\deriv}[2]{\frac{\dd #1}{\dd #2}}
\newcommand{\pderiv}[2]{\frac{\partial #1}{\partial #2}}
\newcommand{\abs}[1]{\left| #1 \right|}
\newcommand{\norm}[1]{\left\| #1 \right\|}
